\newpage
\section{Implementierung der Evaluationsumgebung}

\subsection{Demo-WebGIS-Anwendung}

Die Demo-WebGIS-Anwendung stellt eine realistische, aber kontrollierte Testumgebung für die Evaluation bereit. Sie ist als Single-Page-Webanwendung umgesetzt und dient der Erkundung von weltweiten Wetterdaten und -vorhersagen. Die Anwendung ist unter \url{https://nils-gb156.github.io/ba-webgis-llm-e2e/} erreichbar und kann auch lokal gestartet werden. Der Quellcode liegt im Repository der Arbeit (siehe Fußnote~\ref{fn:repo}).

Abbildung~\ref{fig:demo_webgis_anwendung} zeigt die Oberfläche der Anwendung. Den Mittelpunkt bildet die Karte. Auf der linken Seite befinden sich die Layersteuerung und die Legende, die über Buttons in der Werkzeugleiste am unteren Rand optional ein- und ausgeblendet werden können. Die Layersteuerung ermöglicht das Umschalten zwischen Hintergrundkarten sowie das Ein- und Ausblenden von Layern. Die Legende darunter zeigt die Symbolerklärungen der jeweils aktiven Layer.

\begin{figure}[H]
    \centering
    \includegraphics[width=\textwidth]{demo_webgis_anwendung}
    \caption[Demo-WebGIS-Anwendung]{Oberfläche der Demo-WebGIS-Anwendung nach einem Klick in die Karte. Links Layersteuerung und Legende, oben mittig die Freitextsuche des Geocoders, rechts das Infopanel mit der Wettervorhersage der geklickten Koordinate und Feature-Informationen.}
    \label{fig:demo_webgis_anwendung}
\end{figure}

Die Anwendung bietet drei verschiedene Hintergrundkarten: Carto Light (Standard), Carto Dark und OpenStreetMap. Initial sichtbar sind zwei Punktlayer des Deutschen Wetterdienstes (\acs{DWD})\acused{DWD} für EUCOS-Bodenstationen und UV-Index-Stationen. Zusätzlich ist noch ein flächendeckender Layer der aktuellen Temperatur von OpenWeatherMap sichtbar. Über die Layersteuerung können weitere flächendeckende Layer für Niederschlag und Wolkenbedeckung, ebenfalls von OpenWeatherMap, sowie ein europaweiter UV-Index-Layer des \ac{DWD} eingeblendet werden.

Ein Klick auf die Karte öffnet ein Infopanel, das den Ort und die Wettervorhersage für die nächsten drei Tage im Drei-Stunden-Takt anzeigt. Die insgesamt 24 Einträge werden über die OpenWeatherMap-\ac{API} abgerufen und enthalten jeweils eine kurze Beschreibung des Wetters, Temperatur, Luftfeuchtigkeit sowie Windrichtung und -geschwindigkeit. Wird zusätzlich auf eine EUCOS-Bodenstation oder UV-Index-Station geklickt, werden die Stationsinformationen unterhalb der Wettervorhersage angezeigt. Die EUCOS-Bodenstationen sind als \ac{WFS} eingebunden, sodass die Informationen direkt aus den geladenen Vektor-Features stammen. Die UV-Index-Stationen werden hingegen als \ac{WMS} bereitgestellt, und die Informationen werden über eine GetFeatureInfo-Anfrage abgerufen.

Oben in der Mitte der Karte befindet sich ein Feld für eine Freitextsuche, die über die Nominatim-Geocoding-\ac{API} nach Orten sucht. Bei Auswahl eines Suchergebnisses wird die Karte auf den Ort zentriert und das Infopanel mit der Wettervorhersage angezeigt.

Über die Werkzeugleiste lässt sich außerdem zur initialen Ausdehnung zurücknavigieren und ebenso kann der Kartenausschnitt vergrößert oder verkleinert werden. Zwei weitere Buttons rechts daneben öffnen verschiebbare Floating-Panels für Mess- und Druckfunktion. Das Messwerkzeug ermöglicht das direkte Messen von Linien und Flächen in der Karte. Mit der Druckfunktion kann der aktuelle Kartenausschnitt mit allen aktiven Layern als PDF oder PNG exportiert werden. Im Footer werden die aktuellen Kartenkoordinaten in WGS84 (World Geodetic System 1984) sowie eine Maßstabsleiste und der Maßstab angezeigt.

Die Anwendung wurde mit Open Pioneer Trails entwickelt und basiert dementsprechend auf OpenLayers, React, Chakra UI und TypeScript. Sie nutzt die verfügbaren \ac{OPT}-Pakete für Layersteuerung, Kartennavigation, Legende, Messung, Druck und Koordinatenanzeige. Für die Umsetzung dieser Demo-Anwendung eignet sich \ac{OPT} durch die freie Gestaltbarkeit von Layout und Funktionalität bei gleichzeitig vorhandenen \ac{GIS}-Widgets für genau diese Funktionen. Die Kartenprojektion ist \ac{EPSG}:3857, die Koordinatenanzeige im Footer transformiert diese zu \ac{EPSG}:4326.

\subsection{Testbarkeitsmaßnahmen der Anwendung}

Die Anwendung wurde so umgesetzt, dass ihr Zustand für einen \ac{E2E}-Test prüfbar ist. Auf \ac{DOM}-Ebene geschieht das über \code|data-testid|-Attribute nach der in Kapitel~2.3.2 beschriebenen Trails-Konvention. Insgesamt sind 39 Attribute vergeben, von denen 24 bereits im initialen Seitenzustand sichtbar sind. Die restlichen erscheinen erst nach einer Interaktion.

Für den Kartenzustand reicht das nicht aus, da die Karte auf einem Canvas gerendert wird (siehe Kapitel~2.4). Aktive Layer, Zoomstufe, Zentrum und Highlight sind über das \ac{DOM} nicht direkt auslesbar. Die Anwendung exponiert deshalb das Kartenmodell als globales JavaScript-Objekt \texttt{\_\_openPioneerMap}. Die Zuweisung erfolgt in der Kartenkomponente in einem \code|useEffect|, sobald das Objekt verfügbar ist. Exponiert wird dabei nicht die OpenLayers-Karte selbst, sondern das MapModel von \ac{OPT}. Über dessen Layer-Sammlung (\code|map.layers|) lassen sich beispielsweise die aktive Hintergrundkarte und die Sichtbarkeiten der Layer erreichen. Auf Zoomstufe, Zentrum und Highlight kann das MapModel über \code|map.olMap| zugreifen.

Die Datei \texttt{map-model-helpers.ts} stellt fertige Funktionen für den Zugriff auf dieses Objekt bereit. Jede Funktion enthält einen \code|page.evaluate()|-Aufruf, liest den Zustand im Browser-Kontext aus und gibt einen prüfbaren Wert zurück. Sie sind bewusst ausgelagert, da sie unabhängig von dieser Anwendung sind und in anderen \ac{OPT}-Anwendungen wiederverwendet werden können. Passend zu den Use Cases stellt die Datei die Funktionen \code|getActiveBaseLayerTitle()|, \code|isLayerRendered()|, \code|getMapZoomLevel()|, \code|getMapCenter()| und \code|getHighlightedCoordinate()| bereit.

Diese Maßnahmen sind Eigenschaften der Anwendung und keine Bestandteile des \ac{UI}-Kontexts. Sie sind in allen Stufen unverändert vorhanden. Variabel ist nur, ob dem \ac{LLM} im Prompt mitgeteilt wird, dass sie existieren (vgl. Kapitel~4.1).

\subsection{Use-Case-Sammlung}

Für das Experiment wurden zehn Use Cases (\texttt{use\_cases.md}) erstellt, die typische Nutzerinteraktionen mit der Demo-WebGIS-Anwendung abbilden. Sie folgen dem in Kapitel~4.2 beschriebenen Format und verteilen sich auf drei einfache, zwei mittlere und fünf schwere Use Cases. Tabelle~\ref{tab:use-cases} zeigt eine Übersicht, die vollständigen Beschreibungen sind in Anhang~\ref{app:use-cases} hinterlegt.

\begin{table}[htbp]
    \centering
    \renewcommand{\arraystretch}{1.3}
    \begin{tabularx}{\textwidth}{@{}c >{\hsize=1.2\hsize}X c >{\hsize=0.8\hsize}X@{}}
        \toprule
        \textbf{Nr} & \textbf{Titel} & \textbf{Komplexität} & \textbf{Kernfunktion} \\
        \midrule
        01  & Layer Switcher ein-/ausblenden          & easy   & Panel-Sichtbarkeit   \\
        02  & Hintergrundkarte wechseln               & easy   & Basemap-Wechsel      \\
        03  & Zoom-Buttons verwenden                  & easy   & Zoom-Steuerung       \\
        04  & UV-Index-Overlay aktivieren             & medium & Layer-Rendering      \\
        05  & Niederschlags-Overlay + Legende         & medium & Layer + Legende      \\
        06  & Wettervorhersage per Kartenklick        & hard   & Kartenklick-Abfrage  \\
        07  & Stations-Info abfragen (WMS + WFS)      & hard   & GetFeatureInfo, GetFeature \\
        08  & Distanz messen                          & hard   & Messwerkzeug         \\
        09  & Karte als PNG exportieren               & hard   & Druckfunktion        \\
        10 & Layer konfigurieren, Suche + Vorhersage & hard   & Geocoder + Workflow  \\
        \bottomrule
    \end{tabularx}
    \caption[Übersicht der Use Cases mit Komplexität und getesteter Kernfunktion]{Übersicht der zehn Use Cases mit Komplexitätsstufe und der jeweils getesteten Kernfunktion der Demo-Anwendung. Die vollständigen Beschreibungen mit Vorbedingungen, Schritten und erwarteten Ergebnissen stehen in Anhang~\ref{app:use-cases}.}
    \label{tab:use-cases}
\end{table}

Die einfachen Use Cases prüfen einzelne Bedienelemente über den \ac{DOM}, etwa das Ein- und Ausblenden der Layersteuerung, den Wechsel der Hintergrundkarte oder die Bedienung der Zoom-Buttons. Die mittleren Use Cases aktivieren zunächst ausgeblendete Layer und prüfen anschließend einen daraus resultierenden Zustand. Bei UC-04 ist es das tatsächliche Rendern der Kacheln auf der Karte und bei UC-05 die Aktualisierung der Legende.

Die schweren Use Cases hingegen erfordern Karteninteraktionen, asynchrone Prozesse oder das Zusammenspiel mehrerer Komponenten. UC-06 prüft die Wettervorhersage durch einen Klick auf eine beliebige Position auf der Karte, während UC-07 das gezielte Treffen einer vorgegebenen Koordinate erfordert. Hierbei muss das \ac{LLM} die Koordinate zuerst in Pixelkoordinaten umrechnen, bevor der Klick simuliert werden kann. UC-08 (Streckenmessung) verlangt eine mehrstufige Zeichnung auf dem Canvas inklusive Doppelklick zum Abschluss, und UC-09 prüft den PNG-Export der aktuellen Kartenansicht. Das Besondere an diesen beiden Use Cases ist, dass diese Funktionen direkt über \ac{OPT}-Komponenten eingebunden sind. Bei der Erstellung der WebGIS-Anwendung lässt sich daher nicht beeinflussen, wie im Test mit diesen Komponenten interagiert werden kann. Der komplexeste ist UC-10, da er Layer-Konfiguration, Geocoder-Suche, Kartennavigation und Wettervorhersage in einem einzigen Ablauf kombiniert.

\subsection{Manuell erstellte Referenztests}

Für jeden Use Case wurde manuell ein Playwright-Test erstellt. Die Erstellung erfolgte vor der \ac{LLM}-Generierung, um eine Beeinflussung zu vermeiden. Die Referenztests dienen zum einen als fachlich korrekte Referenzimplementierung, die zeigt, wie eine saubere Umsetzung eines Use Cases aussieht, und zum anderen bilden sie den Vergleichsmaßstab für die LLM-generierten Tests in der Evaluation.

Die Tests sind mit Playwright in TypeScript umgesetzt. Pro Use Case gibt es eine eigene Spec-Datei nach dem Namensschema \texttt{uc-01-*.spec.ts} bis \texttt{uc-10-*.spec.ts}. Jeder Test folgt demselben Aufbau, der die Struktur des Use Cases widerspiegelt: Zuerst wird geprüft, ob die Vorbedingungen erfüllt sind, dann werden die Schritte ausgeführt und zuletzt die erwarteten Ergebnisse überprüft. Listing~\ref{lst:uc-04-uv-index-overlay} zeigt diesen Aufbau am Beispiel von UC-04.

\begin{listing}[H]
\caption{UC-04: UV-Index-Layer einblenden und prüfen, ob er auf der Karte angezeigt wird (Playwright-Test)}
\label{lst:uc-04-uv-index-overlay}
\inputminted{typescript}{Quellcode/uc-04-uv-index-overlay.spec.ts}
\end{listing}

Bevor ein Test mit der Webanwendung interagiert, muss sichergestellt sein, dass die Anwendung vollständig geladen ist. Statt einer festen Wartezeit (Timeout), die unzuverlässig ist, wird mit \code|page.waitForFunction(() => globalThis.__openPioneerMap != null)| gewartet, bis das Kartenmodell verfügbar ist. Sobald das Objekt existiert, ist die Karte einsatzbereit. Der Test wartet damit auf ein deterministisches Signal und nicht auf eine geschätzte Dauer.

Die Locator-Strategie basiert auf der in Kapitel~2.2.2 beschriebenen Hierarchie und nutzt bevorzugt \code|getByTestId|, dann \code|getByRole| und dann \code|getByText|. Die Selektoren müssen zudem bewusst präzise gewählt werden, da es sonst zu Fehlzugriffen kommen kann. Zum Beispiel wird in UC-04 der UV-Index-Layer über \code|getByRole("checkbox", { name: "UV-Index", exact: true })| angesprochen. Die Option \code|exact| ist hier notwendig, weil sonst auch der Layer \enquote{UV-Index Stations} mitgetroffen werden würde, was im Rahmen des Use Cases falsch wäre. Eine Ausnahme ist UC-09. Dort werden zwei Elemente über \ac{CSS}-Klassen angesprochen, weil die betreffenden internen Elemente der Druckfunktion von Open Pioneer Trails keine eigene Test-ID besitzen.

Klicks auf der Karte erfolgen über \code|page.mouse.click(x, y)| mit Pixelkoordinaten, die aus der \code|boundingBox()| des Kartencontainers berechnet werden. Bei UC-06 genügt ein Klick auf eine beliebige Position, z.\,B. die Kartenmitte. Bei UC-07 muss dagegen ein konkreter Punkt getroffen werden. Dafür wird die bekannte Koordinate (\ac{EPSG}:3857) über die OpenLayers-Funktion \code|getPixelFromCoordinate()| in eine Canvas-Pixelposition umgerechnet.

In UC-04 lässt sich das Aktivieren des UV-Index-Layers auf \ac{DOM}-Ebene über den Zustand des Toggles prüfen (\code|toBeChecked|). Dass der Layer in der Layersteuerung aktiviert ist, bedeutet jedoch nicht, dass die Umschaltung im Kartenmodell angekommen ist. Da der Layer asynchron geladen wird, prüft der Test dies zusätzlich wiederholt mit \code|await expect.poll(() => isLayerRendered(page, UVI_LAYER_TITLE)).toBe(true)|.
